% GENERATED FILE. Edit canonical lessons, then run npm run generate:books.
% canonical-source-hash: fnv1a64:084228d5fd7a594c
% canonical-lessons: GU-C182-ajkal, GU-C182-kayam, GU-C182-acanak, GU-C182-yad-rakhvu, GU-C182-gumavvu, GU-R182-review-177-182, GU-R182-review-final-182

\chapter{These days and for good}
\label{ch:gu-these-days-and-for-good}
\hlchaptermodality{182}

\begin{chapteropening}
\textbf{By the end of this chapter:} \emph{I can talk about these days, what will last and what I will keep in mind.}

Complete the last lesson of chapter 182: I can talk about these days, what will last and what I will keep in mind.
\end{chapteropening}

\section[\texorpdfstring{\gu{આજકાલ} (ājkāl)}{આજકાલ (ājkāl)}]{\gu{આજકાલ} (ājkāl) — these days, nowadays}

\label{lesson:GU-C182-ajkal}



\begin{quote}
Before the new one: say the Gujarati for to cancel, then the Gujarati for to extend.
\end{quote}

\subsection*{You'll want to know: \gu{આજકાલ}}
\textbf{\gu{આજકાલ}} — \emph{ājkāl} — \textquotedblleft{}these days, nowadays\textquotedblright{}.

\begin{quote}
\textbf{\gu{આજકાલ} \gu{હું} \gu{વહેલો} \gu{ઉઠું} \gu{છું}.} — \emph{ājkāl hũ vahelo uṭhũ chũ} — \textquotedblleft{}These days I get up early.\textquotedblright{}
\end{quote}

\begin{cousinweb}
\textbf{\gu{આજકાલ}} (\emph{ājkāl}) is built from two taught words, \emph{āj}, \textquotedblleft{}today\textquotedblright{}, and \emph{kāl}, \textquotedblleft{}tomorrow, yesterday\textquotedblright{}, so it means about now.
\end{cousinweb}

\begin{grammarlens}[title={What you've built}]
One more word for everyday life at A2.
\end{grammarlens}

\subsection*{Your turn}
\begin{itemize}
\raggedright
  \item \emph{Say it:} \emph{ājkāl}
  \item \emph{Say it:} \emph{ājkāl}, once more
  \item \emph{Say it:} \emph{ājkāl hũ vahelo uṭhũ chũ} — These days I get up early.
  \item \emph{Recall:} say the Gujarati for to cancel, then the Gujarati for to extend, then say \emph{ājkāl} again
\end{itemize}

\begin{tcolorbox}[breakable,colback=teal!4,colframe=teal!35!black,title={Before you move on}]
What does \gu{આજકાલ} mean? (\textquotedblleft{}These days, nowadays\textquotedblright{}.) Say it once more.
\end{tcolorbox}
\section[\texorpdfstring{\gu{કાયમ} (kāyam)}{કાયમ (kāyam)}]{\gu{કાયમ} (kāyam) — permanently, for good}

\label{lesson:GU-C182-kayam}



\begin{quote}
Before the new one: say the Gujarati for to extend, then the Gujarati for these days.
\end{quote}

\subsection*{You'll want to know: \gu{કાયમ}}
\textbf{\gu{કાયમ}} — \emph{kāyam} — \textquotedblleft{}permanently, for good\textquotedblright{}.

\begin{quote}
\textbf{\gu{આ} \gu{ઘર} \gu{કાયમ} \gu{અમારું} \gu{રહેશે}.} — \emph{ā ghar kāyam amārũ raheśe} — \textquotedblleft{}This house will be ours for good.\textquotedblright{}
\end{quote}

\begin{cousinweb}
\textbf{\gu{કાયમ}} (\emph{kāyam}) is a loan from Arabic \emph{qāʾim}, \textquotedblleft{}standing, lasting\textquotedblright{}, which reached Gujarati through Persian. Its Arabic root is \emph{q-w-m}, \textquotedblleft{}to stand\textquotedblright{} (Lane).
\end{cousinweb}

\begin{grammarlens}[title={What you've built}]
One more word for everyday life at A2.
\end{grammarlens}

\subsection*{Your turn}
\begin{itemize}
\raggedright
  \item \emph{Say it:} \emph{kāyam}
  \item \emph{Say it:} \emph{kāyam}, once more
  \item \emph{Say it:} \emph{ā ghar kāyam amārũ raheśe} — This house will be ours for good.
  \item \emph{Recall:} say the Gujarati for to extend, then the Gujarati for these days, then say \emph{kāyam} again
\end{itemize}

\begin{tcolorbox}[breakable,colback=teal!4,colframe=teal!35!black,title={Before you move on}]
What does \gu{કાયમ} mean? (\textquotedblleft{}Permanently, for good\textquotedblright{}.) Say it once more.
\end{tcolorbox}
\section[\texorpdfstring{\gu{અચાનક} (acānak)}{અચાનક (acānak)}]{\gu{અચાનક} (acānak) — suddenly}

\label{lesson:GU-C182-acanak}



\begin{quote}
Before the new one: say the Gujarati for these days, then the Gujarati for permanently.
\end{quote}

\subsection*{You'll want to know: \gu{અચાનક}}
\textbf{\gu{અચાનક}} — \emph{acānak} — \textquotedblleft{}suddenly\textquotedblright{}.

\begin{quote}
\textbf{\gu{કદાચ} \gu{તે} \gu{અચાનક} \gu{આવશે}.} — \emph{kadāc te acānak āvśe} — \textquotedblleft{}Perhaps he will turn up suddenly.\textquotedblright{}
\end{quote}

\begin{cousinweb}
The origin of \textbf{\gu{અચાનક}} (\emph{acānak}) is not settled. EtymDB lists it as borrowed from Hindi \emph{acānak}, and CDIAL lists no Gujarati form.
\end{cousinweb}

\begin{grammarlens}[title={What you've built}]
One more word for everyday life at A2.
\end{grammarlens}

\subsection*{Your turn}
\begin{itemize}
\raggedright
  \item \emph{Say it:} \emph{acānak}
  \item \emph{Say it:} \emph{acānak}, once more
  \item \emph{Say it:} \emph{kadāc te acānak āvśe} — Perhaps he will turn up suddenly.
  \item \emph{Recall:} say the Gujarati for these days, then the Gujarati for permanently, then say \emph{acānak} again
\end{itemize}

\begin{tcolorbox}[breakable,colback=teal!4,colframe=teal!35!black,title={Before you move on}]
What does \gu{અચાનક} mean? (\textquotedblleft{}Suddenly\textquotedblright{}.) Say it once more.
\end{tcolorbox}
\section[\texorpdfstring{\gu{યાદ} \gu{રાખવું} (yād rākhvũ)}{યાદ રાખવું (yād rākhvũ)}]{\gu{યાદ} \gu{રાખવું} (yād rākhvũ) — to keep in mind}

\label{lesson:GU-C182-yad-rakhvu}



\begin{quote}
Before the new one: say the Gujarati for permanently, then the Gujarati for suddenly.
\end{quote}

\subsection*{You'll want to know: \gu{યાદ} \gu{રાખવું}}
\textbf{\gu{યાદ} \gu{રાખવું}} — \emph{yād rākhvũ} — \textquotedblleft{}to keep in mind\textquotedblright{}.

\begin{quote}
\textbf{\gu{હું} \gu{આ} \gu{તારીખ} \gu{યાદ} \gu{રાખીશ}.} — \emph{hũ ā tārīkh yād rākhīś} — \textquotedblleft{}I will remember this date.\textquotedblright{}
\end{quote}

Word for word, to keep in memory.

\begin{cousinweb}
\textbf{\gu{યાદ} \gu{રાખવું}} (\emph{yād rākhvũ}) is built from \emph{yād}, \textquotedblleft{}memory\textquotedblright{}, usually traced to Persian \emph{yād}, \textquotedblleft{}memory\textquotedblright{}, and the inherited \emph{rākhvũ}, \textquotedblleft{}to keep\textquotedblright{}, which CDIAL traces to Sanskrit \emph{rakṣati}, \textquotedblleft{}guards\textquotedblright{} (CDIAL 10547).
\end{cousinweb}

\begin{grammarlens}[title={What you've built}]
One more word for everyday life at A2.
\end{grammarlens}

\subsection*{Your turn}
\begin{itemize}
\raggedright
  \item \emph{Say it:} \emph{yād rākhvũ}
  \item \emph{Say it:} \emph{yād rākhvũ}, once more
  \item \emph{Say it:} \emph{hũ ā tārīkh yād rākhīś} — I will remember this date.
  \item \emph{Recall:} say the Gujarati for permanently, then the Gujarati for suddenly, then say \emph{yād rākhvũ} again
\end{itemize}

\begin{tcolorbox}[breakable,colback=teal!4,colframe=teal!35!black,title={Before you move on}]
What does \gu{યાદ} \gu{રાખવું} mean? (\textquotedblleft{}To keep in mind\textquotedblright{}.) Say it once more.
\end{tcolorbox}
\section[\texorpdfstring{\gu{ગુમાવવું} (gumāvvũ)}{ગુમાવવું (gumāvvũ)}]{\gu{ગુમાવવું} (gumāvvũ) — to let slip, to lose}

\label{lesson:GU-C182-gumavvu}



\begin{quote}
Before the new one: say the Gujarati for suddenly, then the Gujarati for to keep in mind.
\end{quote}

\subsection*{You'll want to know: \gu{ગુમાવવું}}
\textbf{\gu{ગુમાવવું}} — \emph{gumāvvũ} — \textquotedblleft{}to let slip, to lose\textquotedblright{}.

\begin{quote}
\textbf{\gu{તમે} \gu{સમય} \gu{ગુમાવશો}.} — \emph{tame samay gumāvśo} — \textquotedblleft{}You will lose time.\textquotedblright{}
\end{quote}

\begin{cousinweb}
\textbf{\gu{ગુમાવવું}} (\emph{gumāvvũ}) is a tadbhava. CDIAL lists it, glossed \textquotedblleft{}to squander, lose, pass (time)\textquotedblright{}, under Sanskrit \emph{gamayati}, \textquotedblleft{}causes to go\textquotedblright{} (CDIAL 4028).
\end{cousinweb}

\begin{grammarlens}[title={What you've built}]
One more word for everyday life at A2.
\end{grammarlens}

\subsection*{Your turn}
\begin{itemize}
\raggedright
  \item \emph{Say it:} \emph{gumāvvũ}
  \item \emph{Say it:} \emph{gumāvvũ}, once more
  \item \emph{Say it:} \emph{tame samay gumāvśo} — You will lose time.
  \item \emph{Recall:} say the Gujarati for suddenly, then the Gujarati for to keep in mind, then say \emph{gumāvvũ} again
\end{itemize}

\begin{tcolorbox}[breakable,colback=teal!4,colframe=teal!35!black,title={Before you move on}]
What does \gu{ગુમાવવું} mean? (\textquotedblleft{}To let slip, to lose\textquotedblright{}.) Say it once more.
\end{tcolorbox}
\section[(dialogue)]{Time and plans words, a review}

\label{lesson:GU-R182-review-177-182}



\begin{quote}
Say the Gujarati for again and again and now and then.
\end{quote}

\subsection*{Your turn}
Say these aloud:

\begin{itemize}
\raggedright
  \item again and again, now and then, occasion, regular, everyday
  \item to start, to finish, on time, remaining, to go by
  \item plan, to decide, programme, beforehand, certain
  \item visit, invitation, to wait, to stay on, to pass
  \item to hurry, to postpone, to avoid, to cancel, to extend
  \item these days, permanently, suddenly, to keep in mind, to let slip
\end{itemize}

\begin{tcolorbox}[breakable,colback=teal!4,colframe=teal!35!black,title={Before you move on}]
Again and again? (\textbf{\gu{વારંવાર}}.) Now and then? (\textbf{\gu{અવારનવાર}}.) Occasion? (\textbf{\gu{વખત}}.) Regular? (\textbf{\gu{નિયમિત}}.) Everyday? (\textbf{\gu{રોજિંદું}}.) To start? (\textbf{\gu{શરૂ} \gu{કરવું}}.) To finish? (\textbf{\gu{પતાવવું}}.) On time? (\textbf{\gu{સમયસર}}.) Remaining? (\textbf{\gu{બાકી}}.) To go by? (\textbf{\gu{વીતવું}}.) Plan? (\textbf{\gu{યોજના}}.) To decide? (\textbf{\gu{નક્કી} \gu{કરવું}}.) Programme? (\textbf{\gu{કાર્યક્રમ}}.) Beforehand? (\textbf{\gu{પહેલેથી}}.) Certain? (\textbf{\gu{ચોક્કસ}}.) Visit? (\textbf{\gu{મુલાકાત}}.) Invitation? (\textbf{\gu{આમંત્રણ}}.) To wait? (\textbf{\gu{રાહ} \gu{જોવી}}.) To stay on? (\textbf{\gu{રોકાવું}}.) To pass? (\textbf{\gu{વિતાવવું}}.) To hurry? (\textbf{\gu{ઉતાવળ} \gu{કરવી}}.) To postpone? (\textbf{\gu{મુલતવી} \gu{રાખવું}}.) To avoid? (\textbf{\gu{ટાળવું}}.) To cancel? (\textbf{\gu{રદ} \gu{કરવું}}.) To extend? (\textbf{\gu{લંબાવવું}}.) These days? (\textbf{\gu{આજકાલ}}.) Permanently? (\textbf{\gu{કાયમ}}.) Suddenly? (\textbf{\gu{અચાનક}}.) To keep in mind? (\textbf{\gu{યાદ} \gu{રાખવું}}.) To let slip? (\textbf{\gu{ગુમાવવું}}.)
\end{tcolorbox}
\section[(dialogue)]{These days and for good — once more}

\label{lesson:GU-R182-review-final-182}



\begin{quote}
Say the Gujarati for these days and permanently.
\end{quote}

\subsection*{Your turn}
\begin{itemize}
\raggedright
  \item \emph{Say it:} these days, permanently, suddenly, to keep in mind, to let slip
\end{itemize}

\begin{tcolorbox}[breakable,colback=teal!4,colframe=teal!35!black,title={Before you move on}]
These days? (\textbf{\gu{આજકાલ}}.) Permanently? (\textbf{\gu{કાયમ}}.) Suddenly? (\textbf{\gu{અચાનક}}.) To keep in mind? (\textbf{\gu{યાદ} \gu{રાખવું}}.) To let slip? (\textbf{\gu{ગુમાવવું}}.)
\end{tcolorbox}
